\documentclass[10pt,a4paper]{article}

% ---------- Packages ----------
\usepackage[utf8]{inputenc}
\usepackage[T1]{fontenc}
\usepackage{lmodern}
\usepackage[margin=1.8cm,top=1.5cm,bottom=1.5cm]{geometry}
\usepackage{enumitem}
\usepackage{titlesec}
\usepackage{xcolor}
\usepackage[hidelinks]{hyperref}

\ifdefined\pdfgentounicode
  \input{glyphtounicode}
  \pdfgentounicode=1
\fi

\hypersetup{
  pdftitle={Adrian René López Cintas - CV},
  pdfauthor={Adrian René López Cintas},
  pdflang={en}
}

\definecolor{accent}{HTML}{1F3A5F}
\definecolor{muted}{HTML}{555555}

\pagestyle{empty}
\setlength{\parindent}{0pt}
\setlength{\parskip}{0pt}
\linespread{1.04}

% ---------- Section style ----------
\titleformat{\section}
  {\large\bfseries\color{accent}}
  {}{0pt}
  {}
  [\vspace{-0.6em}\color{accent}\rule{\linewidth}{0.5pt}]
\titlespacing*{\section}{0pt}{1em}{0.5em}

% ---------- Lists ----------
\setlist[itemize]{leftmargin=1.2em,itemsep=1.5pt,topsep=2pt,parsep=0pt,partopsep=0pt}

% ---------- Helpers ----------
% \job{title}{dates}{subtitle}
\newcommand{\job}[3]{%
  \vspace{0.55em}
  {\bfseries #1}\hfill{\small\color{muted}#2}\par
  \ifx\relax#3\relax\else{\small\itshape\color{muted}#3}\par\fi
  \vspace{0.15em}
}
\newcommand{\subhead}[1]{\vspace{0.3em}{\small\bfseries\color{accent}#1}\par\vspace{0.1em}}

\begin{document}

% ---------- Header ----------
\begin{center}
  {\Huge\bfseries Adrian René López Cintas}\\[0.35em]
  {\large\color{accent}Fullstack Engineer \textbullet{} Go \textbullet{} React \textbullet{} SQL}\\[0.35em]
  {\small Havana, Cuba}\\[0.2em]
  {\small \href{mailto:xeraphdev@gmail.com}{xeraphdev@gmail.com} \textbullet{} \href{https://github.com/xeraph-dev}{github.com/xeraph-dev} \textbullet{} \href{https://www.linkedin.com/in/adrian-l%C3%B3pez/}{linkedin.com/in/adrian-lópez}}
\end{center}

% ---------- Summary ----------
\section*{Summary}
Self-taught fullstack engineer with 3 years of Go and SQL on the back end and 6 years of React on the front end, with experience in query performance optimization on large datasets, accessible web interfaces, and migrating legacy applications to React. Currently in the second year of a BSc in Mathematics.

% ---------- Experience ----------
\section*{Experience}
\job{Web Developer (Go / React), Divertimento}{Nov 2022--Nov 2025 \textbullet{} Remote}{}
\subhead{Healthcare care-management platform}
\begin{itemize}
  \item Developed a Go microservice acting as a backend-for-frontend (gateway) that served the platform's web application exclusively, loading large volumes of data from another microservice.
  \item Built the back-end logic of a task-generation system: staff define patient criteria, the system generates follow-up lists, and a detail view lets them contact patients and record outcomes.
  \item Implemented caching strategies in the service's own PostgreSQL database to avoid repeatedly loading large datasets from the upstream service.
  \item Fixed a production report page that timed out on every request: it joined 10 tables over millions of rows against a 30 s limit. Redesigned the queries with CTEs and tighter joins, plus pagination and precise filters, bringing response times to 5–20 s depending on the filter volume. Before, it was slow regardless of how few results it returned.
  \item Designed a reusable Go abstraction for dynamic filtering and pagination over large tables.
  \item Reduced load in many other queries by narrowing the working set early instead of expanding rows and filtering afterwards.
\end{itemize}
\subhead{Other Projects}
\begin{itemize}
  \item Built a full portal from scratch in Astro, implementing a provided design.
  \item Helped build a product launch landing page in Laravel, nearly from scratch.
  \item Worked on the front end of migrating a legacy .NET application to React.
  \item Migrated an end-to-end test suite from Cucumber to Cypress.
\end{itemize}
\job{Front-end Developer (React), Freelance}{Sep 2022--Nov 2022 \textbullet{} Remote}{}
\begin{itemize}
  \item Built a web page for designing databases without needing to know SQL.
  \item Wrote Cypress end-to-end tests for every new feature and bug fix, so that fixed issues could not regress.
  \item Improved site performance with React best practices; contributed to core product development.
\end{itemize}
\job{Web Developer (PHP / React), Divertimento}{Mar 2022--Aug 2022 \textbullet{} Remote}{}
\subhead{School e-learning platform}
\begin{itemize}
  \item Built core features of an online course platform developed for a single school.
  \item Worked under strict WCAG AA accessibility requirements, testing with screen readers on Windows and VoiceOver on macOS to reproduce and fix assistive-technology issues.
  \item Implemented an in-house text-highlighting library from scratch as part of the project's own component set.
  \item Integrated with Java/Spring Boot microservices, occasionally modifying their code.
\end{itemize}
\job{Front-end Developer (React), IJP Software}{Jul 2021--Jul 2022 \textbullet{} Remote \textbullet{} Part-time at first}{}
\begin{itemize}
  \item Worked on a platform built around dynamically generated forms: forms arrived as structured data and the front end rebuilt the UI from them at runtime.
  \item Built the front end of an e-commerce site in Vue.
  \item Helped build a complex product almost from scratch (React, Next.js, Redux, Docker).
  \item Built a mobile app with the core team in a hackathon, working on both the front end (React Native) and the back end (Express).
  \item Contributed to a React Native music-purchasing app.
\end{itemize}
\job{Front-end Developer (React), Freelance}{Jul 2020--Jul 2021 \textbullet{} Remote \textbullet{} Part-time}{}
\begin{itemize}
  \item Built small JavaScript and React projects for clients.
  \item Wrote a script to add watermarks to videos, which led to a client getting back in touch two years later.
\end{itemize}
\job{Drupal Developer, Xetid}{Oct 2019--Aug 2021 \textbullet{} Havana, Cuba}{}
\begin{itemize}
  \item Built two government websites from scratch (Drupal, PHP, PostgreSQL).
  \item Developed a Drupal module for user authentication through EnZona, a Cuban digital payments platform.
  \item Maintained and extended existing government websites by adding new features.
\end{itemize}

% ---------- Projects ----------
\section*{Projects}
\begin{itemize}
  \item \textbf{Guardias}: cross-platform desktop scheduling app in Tcl/Tk, built as a single static binary; in active development. \href{https://github.com/xeraph-dev/guardias}{github.com/xeraph-dev/guardias}
  \item \textbf{Sudoku}: desktop sudoku game built with Svelte and Tauri. \href{https://github.com/xeraph-dev/sudoku}{github.com/xeraph-dev/sudoku}
  \item \textbf{text-highlighter}: text-highlighting library written from scratch for the e-learning platform. \href{https://github.com/xeraph-dev/text-highlighter}{github.com/xeraph-dev/text-highlighter}
  \item \textbf{Channels Watcher Bot}: Telegram bot (Python, Prisma, Docker) that watches channels and forwards messages matching each user's filters. \href{https://github.com/xeraph-dev/channels-watcher-bot}{github.com/xeraph-dev/channels-watcher-bot}
  \item \textbf{SRI}: information retrieval system for documents based on the vector space model (Python), built in my first year at university. \href{https://github.com/xeraph-dev/sri}{github.com/xeraph-dev/sri}
  \item \textbf{Space Invaders}: game written in Racket as the final project of How to Design Programs. \href{https://github.com/xeraph-dev/space-invaders.rkt}{github.com/xeraph-dev/space-invaders.rkt}
  \item \textbf{Calendar Generator}: calendar generator in Swift that outputs calendars in several formats. \href{https://github.com/xeraph-dev/calendar-generator}{github.com/xeraph-dev/calendar-generator}
  \item \textbf{VSCode Profile Manager}: macOS app in Swift for managing Visual Studio Code profiles. \href{https://github.com/xeraph-dev/vscode-profile-manager-macos}{github.com/xeraph-dev/vscode-profile-manager-macos}
  \item \textbf{Rick and Morty}: macOS client for the Rick and Morty API, written in Swift. \href{https://github.com/xeraph-dev/rick-and-morty-macos}{github.com/xeraph-dev/rick-and-morty-macos}
  \item \textbf{BettyCook}: mobile recipe-book app built with Flutter. \href{https://github.com/xeraph-dev/bettycook}{github.com/xeraph-dev/bettycook}
\end{itemize}

% ---------- Education ----------
\section*{Education}
\textbf{BSc in Mathematics} (in progress), University of Havana \hfill {\small\color{muted}Sep 2025--Present}\par

% ---------- Skills ----------
\section*{Skills}
\begin{itemize}
  \item[] \textbf{Back-end:} Go, SQL (PostgreSQL, MySQL), REST APIs, Node.js/Express, Docker, Linux
  \item[] \textbf{Front-end:} React, React Native, TypeScript, Next.js, Vue, accessibility (WCAG AA), end-to-end testing (Cypress)
  \item[] \textbf{Other:} Tcl/Tk, PHP, Drupal, Git, Haskell, Racket, Python, Swift, Flutter, Java (basic)
  \item[] \textbf{Process:} Scrum, Jira
\end{itemize}

% ---------- Languages ----------
\section*{Languages}
Spanish (native) \textbullet{} English (intermediate)

\end{document}
